% Generated draft: requires mathematical and source review.
\documentclass[11pt]{article}
\usepackage[margin=25mm]{geometry}
\usepackage{fontspec}
\setmainfont{Latin Modern Roman}
\usepackage{amsmath,amssymb,mathtools,unicode-math}
\setmathfont{Latin Modern Math}
\usepackage{enumitem,needspace,xurl,hyperref}
\setlength{\parindent}{0pt}
\setlength{\parskip}{5pt}
\setlength{\emergencystretch}{3em}
\newcommand{\N}{\mathbb{N}}
\newcommand{\Z}{\mathbb{Z}}
\newcommand{\Q}{\mathbb{Q}}
\newcommand{\R}{\mathbb{R}}
\newcommand{\C}{\mathbb{C}}
\newcommand{\F}{\mathbb{F}}
\newcommand{\A}{\mathbb{A}}
\newcommand{\PP}{\mathbb{P}}
\newcommand{\E}{\mathbb{E}}
\newcommand{\eps}{\varepsilon}
\newcommand{\dd}{\,\mathrm{d}}
\begin{document}
\section*{UnsolvedMath ID: 701101}
\textbf{Source problem code:} catalogue entry
\subsection{Definitions and mathematical statement}
A finite presentation of a group is an expression
\[
G\cong\langle s_1,\ldots,s_d\mid r_1,\ldots,r_k\rangle,
\]
where \(d,k\) are finite nonnegative integers and the relators \(r_j\) are words in the generators and their inverses.

For the term \emph{branch group}, use the following standard rooted-tree formulation. Let \(T\) be an infinite rooted tree such that every vertex at distance \(n\) from the root has exactly \(m_n\) children, where \(2\leq m_n<\infty\). Let \(L_n\) denote its level of vertices at distance \(n\) from the root. Suppose that \(G\) acts faithfully on \(T\) by root-preserving automorphisms. For a vertex \(v\), let \(T_v\) be the subtree consisting of \(v\) and all its descendants, and define the rigid stabilizer
\[
\operatorname{rist}_G(v)=\{g\in G: g(w)=w\text{ for every vertex }w\notin T_v\}.
\]
Define the rigid level stabilizer by
\[
\operatorname{rist}_G(n)=\left\langle\operatorname{rist}_G(v):v\in L_n\right\rangle.
\]
The action is a branch action if it is transitive on every \(L_n\) and
\[
[G:\operatorname{rist}_G(n)]<\infty\qquad\text{for every }n\geq0.
\]
A group admitting such an action is called a branch group.

\textbf{Problem.} Does there exist a finitely presented group \(G\) admitting a branch action as defined above?

\textbf{Formulation qualification.} The supplied record does not define ``branch group.'' The rooted-tree definition above makes explicit the conventional interpretation used in this draft; its agreement with the intended source convention requires human review. The supplied registry labels the problem open; that label is not an independently verified assessment.
\subsection{Short English statement}
Does any branch group have a presentation with finitely many generators and finitely many defining relations? Here the branch condition is made explicit through a faithful, level-transitive action on a rooted tree whose rigid level stabilizers all have finite index. The supplied source does not state its definition of branch group.
\subsection{Sources}
\begin{itemize}
\item Local UnsolvedMath registry, permanent ID \(701101\), exact source problem code \texttt{catalogue entry}, title ``Kourovka Notebook Problem 15.13.'' The supplied record gives proposed year \(2002\), status \texttt{open}, and upstream revision \texttt{d371ee917ac9647fbb06cd6fc529deec6f3dda82}. Its creation and update timestamps are both \texttt{2026-10-16T14:00:00Z}; these are supplied registry metadata, not independent source verification.
\item \emph{The Kourovka Notebook (Unsolved Problems in Group Theory)}, 15th issue (2002), Problem 15.13; also identified in the supplied background with the 21st issue edition (2026), arXiv:1401.0300v48. The supplied collection description names E. I. Khukhro and V. D. Mazurov as editors. These bibliographic details are reproduced from the record and have not been independently checked. No external URL is supplied.
\item Attribution supplied: L. Bartholdi, R. I. Grigorchuk, and a third proposer whose name is corrupted. The proposer field ends with the fragment \(\mathrm{Z.}\ \{\backslash v\,S\); the background repeats this fragment and places the detached fragment \(\mathrm{uni}\backslash{}'k\}\) after ``Discussion and literature:''. The third name cannot be transcribed reliably from these fields without correction of the source.
\end{itemize}
\end{document}
